\documentclass[12pt,a4paper]{article}
\usepackage[UTF8]{ctex}
\usepackage{amsmath,amssymb,amsthm}
\usepackage{geometry}

\geometry{margin=2.5cm}

\title{矩阵指数解与稳定性证明}
\author{第11章补充材料}
\date{}

\newtheorem{theorem}{定理}
\newtheorem{lemma}{引理}
\newtheorem{proof}{证明}
\newtheorem{corollary}{推论}

\begin{document}

\maketitle

\section{问题陈述}

考虑向量微分方程：
\begin{equation}
\dot{\mathbf{x}}(t) = A\mathbf{x}(t), \quad \mathbf{x}(0) = \mathbf{x}_0,
\end{equation}
其中$\mathbf{x}(t) \in \mathbb{R}^n$，$A \in \mathbb{R}^{n \times n}$是常数矩阵。

我们要证明：
\begin{enumerate}
\item 方程(1)的解为$\mathbf{x}(t) = e^{At}\mathbf{x}_0$，其中$e^{At}$是矩阵指数。
\item 如果矩阵$A$的所有特征值具有负实部，则解$\mathbf{x}(t)$从任何初始条件$\mathbf{x}_0$收敛到平衡点$\mathbf{0}$。
\end{enumerate}

\section{标量情况的类比}

首先，我们回顾标量一阶微分方程：
\begin{equation}
\dot{x}(t) = ax(t), \quad x(0) = x_0,
\end{equation}
其中$a \in \mathbb{R}$是常数。

通过分离变量法，我们可以求解这个方程：
\begin{align}
\frac{dx}{x} &= a \, dt, \\
\int_{x_0}^{x(t)} \frac{dx}{x} &= \int_0^t a \, d\tau, \\
\ln\left(\frac{x(t)}{x_0}\right) &= at, \\
x(t) &= x_0 e^{at}.
\end{align}

因此，标量微分方程(2)的解为：
\begin{equation}
x(t) = e^{at}x_0.
\end{equation}

\section{矩阵指数的定义}

为了将标量情况推广到向量情况，我们需要定义矩阵指数。

\begin{definition}[矩阵指数]
对于矩阵$A \in \mathbb{R}^{n \times n}$，矩阵指数$e^{At}$定义为：
\begin{equation}
e^{At} = I + At + \frac{(At)^2}{2!} + \frac{(At)^3}{3!} + \cdots = \sum_{k=0}^{\infty} \frac{(At)^k}{k!},
\end{equation}
其中$I$是$n \times n$单位矩阵。
\end{definition}

这个级数对所有矩阵$A$和所有$t \in \mathbb{R}$都收敛。

\section{矩阵指数的性质}

\begin{lemma}
矩阵指数$e^{At}$满足以下性质：
\begin{enumerate}
\item $e^{A \cdot 0} = I$
\item $\frac{d}{dt}e^{At} = Ae^{At} = e^{At}A$
\item $e^{A(t+s)} = e^{At}e^{As}$（半群性质）
\item 如果$AB = BA$，则$e^{(A+B)t} = e^{At}e^{Bt}$
\end{enumerate}
\end{lemma}

\begin{proof}
\begin{enumerate}
\item 直接代入定义(8)，当$t=0$时，所有$k \geq 1$的项都为零，只剩下$I$。

\item 对级数(8)逐项求导：
\begin{align}
\frac{d}{dt}e^{At} &= \frac{d}{dt}\sum_{k=0}^{\infty} \frac{(At)^k}{k!} \\
&= \sum_{k=1}^{\infty} \frac{kA(At)^{k-1}}{k!} \\
&= A\sum_{k=1}^{\infty} \frac{(At)^{k-1}}{(k-1)!} \\
&= A\sum_{j=0}^{\infty} \frac{(At)^j}{j!} = Ae^{At}.
\end{align}
由于矩阵乘法满足结合律，且$A$与$e^{At}$可交换（因为$A$与自己的幂次可交换），所以$Ae^{At} = e^{At}A$。

\item 利用级数展开和Cauchy乘积公式可以证明。

\item 当$AB = BA$时，可以利用二项式定理展开$(A+B)^k$来证明。
\end{enumerate}
\end{proof}

\section{向量微分方程的解}

\begin{theorem}
向量微分方程(1)的解为：
\begin{equation}
\mathbf{x}(t) = e^{At}\mathbf{x}_0.
\end{equation}
\end{theorem}

\begin{proof}
我们需要验证$\mathbf{x}(t) = e^{At}\mathbf{x}_0$满足微分方程(1)和初始条件。

\textbf{验证初始条件：}
根据引理1的性质(1)，$e^{A \cdot 0} = I$，因此：
\begin{equation}
\mathbf{x}(0) = e^{A \cdot 0}\mathbf{x}_0 = I\mathbf{x}_0 = \mathbf{x}_0,
\end{equation}
满足初始条件。

\textbf{验证微分方程：}
对$\mathbf{x}(t) = e^{At}\mathbf{x}_0$关于$t$求导，利用引理1的性质(2)：
\begin{align}
\dot{\mathbf{x}}(t) &= \frac{d}{dt}\left(e^{At}\mathbf{x}_0\right) \\
&= \frac{d}{dt}\left(e^{At}\right)\mathbf{x}_0 \\
&= Ae^{At}\mathbf{x}_0 \\
&= A\mathbf{x}(t),
\end{align}
这正好是微分方程(1)的右侧。因此，$\mathbf{x}(t) = e^{At}\mathbf{x}_0$是方程(1)的解。
\end{proof}

\section{稳定性分析}

现在我们来分析解的稳定性，即何时$\mathbf{x}(t) \to \mathbf{0}$当$t \to \infty$。

\subsection{特征值与特征向量}

设矩阵$A$有$n$个特征值$\lambda_1, \lambda_2, \ldots, \lambda_n$（可能重复），对应的特征向量为$\mathbf{v}_1, \mathbf{v}_2, \ldots, \mathbf{v}_n$。如果$A$可以对角化，则存在可逆矩阵$P$使得：
\begin{equation}
A = P\Lambda P^{-1},
\end{equation}
其中$\Lambda = \text{diag}(\lambda_1, \lambda_2, \ldots, \lambda_n)$是对角矩阵。

\subsection{矩阵指数的特征值分解}

\begin{lemma}
如果$A = P\Lambda P^{-1}$，则：
\begin{equation}
e^{At} = Pe^{\Lambda t}P^{-1} = P\text{diag}(e^{\lambda_1 t}, e^{\lambda_2 t}, \ldots, e^{\lambda_n t})P^{-1}.
\end{equation}
\end{lemma}

\begin{proof}
利用矩阵指数的级数定义：
\begin{align}
e^{At} &= \sum_{k=0}^{\infty} \frac{(At)^k}{k!} \\
&= \sum_{k=0}^{\infty} \frac{(P\Lambda P^{-1}t)^k}{k!} \\
&= \sum_{k=0}^{\infty} \frac{P(\Lambda t)^k P^{-1}}{k!} \\
&= P\left(\sum_{k=0}^{\infty} \frac{(\Lambda t)^k}{k!}\right)P^{-1} \\
&= Pe^{\Lambda t}P^{-1}.
\end{align}

由于$\Lambda$是对角矩阵，$e^{\Lambda t}$也是对角矩阵，其对角线元素为$e^{\lambda_i t}$。
\end{proof}

\subsection{稳定性定理}

\begin{theorem}[稳定性定理]
向量微分方程(1)的解$\mathbf{x}(t) = e^{At}\mathbf{x}_0$从任何初始条件$\mathbf{x}_0$收敛到平衡点$\mathbf{0}$，当且仅当矩阵$A$的所有特征值具有负实部。
\end{theorem}

\begin{proof}
假设$A$可以对角化（对于不能对角化的情况，可以使用Jordan标准型，证明类似）。

根据引理2，解可以写成：
\begin{equation}
\mathbf{x}(t) = e^{At}\mathbf{x}_0 = Pe^{\Lambda t}P^{-1}\mathbf{x}_0 = P\text{diag}(e^{\lambda_1 t}, e^{\lambda_2 t}, \ldots, e^{\lambda_n t})P^{-1}\mathbf{x}_0.
\end{equation}

设$\mathbf{y} = P^{-1}\mathbf{x}_0$，则：
\begin{equation}
\mathbf{x}(t) = P\text{diag}(e^{\lambda_1 t}, e^{\lambda_2 t}, \ldots, e^{\lambda_n t})\mathbf{y}.
\end{equation}

\textbf{充分性（如果所有特征值具有负实部，则解收敛）：}

如果对所有$i$，$\text{Re}(\lambda_i) < 0$，则：
\begin{equation}
|e^{\lambda_i t}| = e^{\text{Re}(\lambda_i) t} \to 0 \quad \text{当} \quad t \to \infty.
\end{equation}

因此，$\text{diag}(e^{\lambda_1 t}, e^{\lambda_2 t}, \ldots, e^{\lambda_n t}) \to \mathbf{0}$（零矩阵），从而：
\begin{equation}
\mathbf{x}(t) = P\text{diag}(e^{\lambda_1 t}, e^{\lambda_2 t}, \ldots, e^{\lambda_n t})\mathbf{y} \to P \cdot \mathbf{0} \cdot \mathbf{y} = \mathbf{0}.
\end{equation}

\textbf{必要性（如果解收敛，则所有特征值具有负实部）：}

如果$\mathbf{x}(t) \to \mathbf{0}$对所有初始条件$\mathbf{x}_0$成立，特别地，对于$\mathbf{x}_0 = P\mathbf{e}_i$（其中$\mathbf{e}_i$是第$i$个标准基向量），我们有：
\begin{equation}
\mathbf{x}(t) = Pe^{\Lambda t}P^{-1}P\mathbf{e}_i = Pe^{\Lambda t}\mathbf{e}_i = P(e^{\lambda_i t}\mathbf{e}_i) = e^{\lambda_i t}P\mathbf{e}_i.
\end{equation}

由于$\mathbf{x}(t) \to \mathbf{0}$且$P\mathbf{e}_i \neq \mathbf{0}$，必须有$e^{\lambda_i t} \to 0$，这意味着$\text{Re}(\lambda_i) < 0$。

因此，对所有$i$，$\text{Re}(\lambda_i) < 0$。
\end{proof}

\section{标量情况的对应}

对于标量微分方程$\dot{x}(t) = ax(t)$：
\begin{itemize}
\item 解为$x(t) = e^{at}x_0$
\item 如果$a < 0$，则$x(t) \to 0$当$t \to \infty$
\item 如果$a > 0$，则$|x(t)| \to \infty$当$t \to \infty$
\item 如果$a = 0$，则$x(t) = x_0$（常数）
\end{itemize}

这与向量情况的结论完全对应：标量$a$对应矩阵$A$的特征值，负的$a$对应具有负实部的特征值。

\section{结论}

我们证明了：
\begin{enumerate}
\item 向量微分方程$\dot{\mathbf{x}}(t) = A\mathbf{x}(t)$的解为$\mathbf{x}(t) = e^{At}\mathbf{x}_0$，这是标量情况$x(t) = e^{at}x_0$的自然推广。

\item 解的稳定性完全由矩阵$A$的特征值决定：
   \begin{itemize}
   \item 如果所有特征值具有负实部，则解从任何初始条件收敛到$\mathbf{0}$（系统是稳定的）。
   \item 如果至少有一个特征值具有正实部，则存在初始条件使得解发散（系统是不稳定的）。
   \item 如果所有特征值具有非正实部，且具有零实部的特征值对应的Jordan块都是一阶的，则解有界但不一定收敛（系统是临界稳定的）。
   \end{itemize}
\end{enumerate}

这些结果构成了线性系统稳定性理论的基础。

\end{document}

